\title{xLSTM: Extended Long Short-Term Memory}

\begin{abstract}
During the 1990s, the Long Short-Term Memory (LSTM) architecture was pioneered through the introduction of gating mechanisms and the constant error carousel. Over subsequent decades, LSTMs have demonstrated remarkable robustness and formed the backbone of many significant deep learning advancements, notably as the foundation for the first Large Language Models (LLMs). Despite these achievements, the emergence of the Transformer framework—centered on parallelizable self-attention—significantly shifted the paradigm by surpassing LSTM-based models at larger scales. In this study, we pose the following fundamental inquiry: What is the potential of LSTMs in language modeling when extended to billions of parameters, utilizing advancements adopted in modern LLMs, yet addressing known shortcomings inherent to LSTMs? To this end, we propose exponential gating, accompanied by rigorous normalization and stabilization strategies. We further innovate on LSTM's memory structure, presenting: (i) sLSTM, which incorporates scalar memory, scalar updates, and revised memory mixing techniques; (ii) mLSTM, characterized by a matrix memory and a covariance-driven update, supporting full parallelization. These LSTM variants are embedded within residual block architectures to form xLSTM blocks, which are consecutively stacked to construct xLSTM models. The combination of exponential gating and novel memory designs empowers xLSTM architectures to rival state-of-the-art Transformer and State Space Models, demonstrating competitive performance and scalability.\\
Code available at: \url{https://github.com/NX-AI/xlstm}
\end{abstract}

\section{Introduction}
\label{sec:intro}

The Long Short-Term Memory (LSTM) ideas~\citep{Hochreiter:91,Hochreiter:97nips,Hochreiter:97}, i.e., the constant error carousel and gating, were introduced to overcome the vanishing gradient problem of recurrent neural networks~\citep{Hochreiter:91,Hochreiter:00book}:

\begin{align}
\label{eq:lstm_idea}
  c_t \ = \ f_t \ c_{t-1} \ + \ 
          i_t \  z_t \ , \quad  
          h_t \ = \ o_t \ \psi({c_t}) \ . 
\end{align}

The cell state $c_{t-1}$ (indicated in green) is incrementally updated by incorporating cell inputs $z_t$, with the degree of update regulated by sigmoid gating mechanisms (depicted in blue). The input gate $i_t$ determines the extent to which new information is introduced, and the forget gate $f_t$ decides how much of the previous cell state is retained. Meanwhile, the output gate $o_t$ modulates the final memory cell output, corresponding to the hidden state $h_t$. Before output, the cell state is either normalized or transformed by $\psi$, after which the output gate produces the hidden state.

LSTMs have achieved notable success across multiple domains \citep{Hochreiter:01,Hochreiter:07,Schmidhuber:15}, dominating text generation tasks up until the introduction of Transformers in 2017~\citep{Vaswani:17}. Their efficacy has been validated in a variety of sequential tasks, such as text synthesis~\citep{Graves:13,Karpathy:15blog}, handwriting creation~\citep{Graves:13}, machine translation between sequences~\citep{Sutskever:14nips}, evaluation of computer code~\citep{Zaremba:14arxiv}, automatic captioning of images~\citep{Karpathy:15,Hossain:19}, source code generation~\citep{Karpathy:15blog}, modeling rainfall and runoff~\citep{Kratzert:18,Kratzert:19}, and hydrological forecasting for flood alerts~\citep{Nearing:24}. In the field of reinforcement learning, LSTMs represent the top-performing architectures for sequential modeling, demonstrated by applications in systems such as AlphaStar for StarCraft II~\citep{Vinyals:17short}, OpenAI Five for Dota 2~\citep{OpenAI:19}, and the magnetic control system for nuclear fusion~\citep{Degrave:22short}. Their capacity for abstraction—skillfully extracting semantic content and maintaining it within their cell state memory~\citep{Karpathy:15blog}—is evidenced by emergent number and syntax units~\citep{Lakretz:19}, linguistic features~\citep{Bau:19}, and sentiment encoding~\citep{Radford:17arxiv}. LSTMs remain integral in high-impact fields~\citep{Degrave:22short,Nearing:24}, proving their robustness and enduring relevance.

Although LSTMs have achieved remarkable performance, they are subject to three primary drawbacks: (i) Inability to revise previously made storage decisions. This shortcoming is illustrated through the {\it Nearest Neighbor Search} problem (refer also to Appendix~\ref{sec:appNearestNeighborSearch}): Given a reference vector, the task requires a sequential scan of a series to identify and return the value attached to its most similar vector at the end. As depicted in the left panel of Figure~\ref{fig:lstmProblems}, the mean squared error indicates that LSTM models are unable to overwrite stored information even if a more suitable vector is encountered later in the sequence; in contrast, our newly introduced xLSTM addresses this via exponential gating. (ii) Constraints on storage capacity, necessitating that information be condensed into scalar cell states. This is demonstrated using {\it Rare Token Prediction}, where, as seen in the right panel of Figure~\ref{fig:lstmProblems}, perplexity results on the Wikitext-103~\citep{Merity:17} dataset are displayed across varying token frequency partitions. Due to its limited capacity, the LSTM is particularly ineffective at predicting infrequent tokens, whereas our xLSTM, featuring a matrix-based memory structure, effectively mitigates this issue. (iii) Restrictions in parallelization caused by the entangling of memories, where recurrent connections between hidden states at successive time steps inherently require sequential computation.

These deficiencies of the LSTM architecture contributed to the rise of Transformers~\citep{Vaswani:17} for language modeling applications. This naturally leads to the question: what level of language modeling performance can be attained by overcoming these issues and extending LSTM architectures to the scale of today’s Large Language Models?

\section{Extended Long Short-Term Memory}
\label{sec:xLSTM}

In order to address the inherent constraints of LSTM architectures, Extended Long Short-Term Memory (xLSTM) incorporates two pivotal modifications to the foundational LSTM mechanism presented in Equation~\eqref{eq:lstm_idea}. These changes—namely, exponential gating mechanisms and innovative memory constructs—expand the LSTM model family by introducing two distinct variants: (i) the sLSTM (refer to Section~\ref{sec:sLSTM}), which employs scalar-valued memory and updates alongside memory mixing, and (ii) the mLSTM (explained in Section~\ref{sec:mLSTM}), distinguished by its use of matrix-based memory together with an update rule based on covariances (outer products), which permits complete parallel computation. Both sLSTM and mLSTM incorporate exponential gating to augment the LSTM's performance. The mLSTM variant omits traditional memory mixing—that is, it discards hidden-to-hidden recurrent connections—to facilitate parallelism. Both variants may be scaled to utilize multiple memory cells, with sLSTM supporting memory interactions among cells. Moreover, sLSTM allows for multiple head configurations, in which inter-head memory mixing is disabled but intra-head, cross-cell mixing persists. By leveraging exponential gating in conjunction with the introduction of multiple heads, sLSTM establishes a novel approach to memory mixing. In contrast, for mLSTM, employing multiple cells and multiple heads yields equivalent functionality.

These new LSTM forms can be integrated within residual block-based modules to create xLSTM blocks (detailed in Section~\ref{sec:xLSTMblock}). Layering such xLSTM blocks via residual connections leads to the formation of complete xLSTM architectures (described in Section~\ref{sec:xLSTMarchitecure}). The organizational structure of the xLSTM and its constituent parts is visualized in Figure~\ref{fig:xlstm_sketch}.

\subsection{Review of the Long Short-Term Memory}
\label{sec:orgLSTM}

The foundational concept of LSTM~\citep{Hochreiter:91,Hochreiter:97nips,Hochreiter:97} incorporated a scalar memory cell, serving as the primary unit for computation and information retention. This design mitigates the issue of vanishing gradients~\citep{Hochreiter:91,Hochreiter:00book} via the constant error carousel, which refers to the way the cell state is updated. The memory cell is structured with three distinct gates: input gate, output gate, and forget gate—the latter was later added by \citet{Gers:00}. The mechanism for updating the LSTM memory cell at each time point $t$ is as follows:

\begin{align}
c_t \ &= \  f_t \ c_{t-1} \ + \ i_t \ z_t &  & &\text{cell state} \\
h_t  \ &= \ o_t \ \tilde{h}_t \ , & \tilde{h}_t \ &= \ \psi \left( c_t\right)
  &\text{hidden state} \\
z_t \ &= \ \varphi \left( \tilde{z}_t \right) \ , 
  &\tilde{z}_t \ &=  \ \mathbf{w}^\top_{z} \ \mathbf{x}_t \ + \
  r_{z}  h_{t-1} \ + \  b_{z} \ \
  &\text{cell input} \\
i_t \ &= \ \sigma \left( \tilde{i}_t  \right) \ , 
  &\tilde{i}_t \ &= \ \mathbf{w}^\top_{i} \ \mathbf{x}_t \ + \
  r_{i}  \ h_{t-1} \ + \  b_{i} \ \
  &\text{input gate} \\
f_t \ &= \ \sigma \left(  \tilde{f}_t \right) \ , 
  &\tilde{f}_t \ &= \ \mathbf{w}^\top_{f} \ \mathbf{x}_t  \ + \
  r_{f}  \ h_{t-1} \ + \  b_{f} \ \
  &\text{forget gate} \\
o_t \ &= \ \sigma \left( \tilde{o}_t \right) \ , 
  &\tilde{o}_t  \ &= \ \mathbf{w}^\top_{o} \ \mathbf{x}_t \ + \
  r_{o}  \ h_{t-1} \ + \  b_{o} \ \
  &\text{output gate} 
\end{align}

The vectors $\mathbf{w}_{z}$, $\mathbf{w}_{i}$, $\mathbf{w}_{f}$, and $\mathbf{w}_{o}$ denote the input weight parameters mapping $\mathbf{x}_t$ to the cell input, input gate, forget gate, and output gate, respectively. The variables $r_{z}$, $r_{i}$, $r_{f}$, and $r_{o}$ serve as the recurrent weight parameters connecting the previous hidden state $h_{t-1}$ to the cell input, input gate, forget gate, and output gate, respectively. The associated bias terms are indicated by $b_{z}$, $b_{i}$, $b_{f}$, and $b_{o}$. The functions $\varphi$ and $\psi$ represent the activation functions for the cell input and hidden state, most often chosen as $\tanh$. The activation function $\psi$ ensures the cell state remains within a bounded range, otherwise it could diverge. All gate activations use the sigmoid function, defined as $\sigma \left( x \right)= 1/(1+\exp(-x))$. In subsequent models, the cell state was structured as a vector of $d$ scalar memory cells, $\mathbf{c}_t \in \mathbb{R}^{d}$, facilitating the introduction of recurrent weight matrices $\mathbf{R} \in \mathbb{R}^{d \times d}$ to blend the outputs of individual memory cells~\citep{Greff:15}; additional discussion can be found in Appendix~\ref{sec:apporgLSTM}. Evidence from ablation analyses supports that each part of the memory cell contributes essentially to overall performance~\citep{Greff:15}.

\subsection{sLSTM}
\label{sec:sLSTM}

In order to enable LSTMs to adapt their storage choices, we incorporate exponential gates (depicted in red), alongside normalization and stabilization mechanisms. Specifically, exponential activation functions are applied to the input and forget gates. For normalization purposes, we add a normalizer state which accumulates the sum of the input gate multiplied by each subsequent forget gate.

The scalar sLSTM forward pass is:

\begin{align}
\label{eq:slstmforward}
c_t \ &= \  f_t \ c_{t-1} \ + \ i_t \ z_t & & 
 &\text{cell state} \\
n_t \ &= \  f_t \ n_{t-1} \ + \ i_t  & & 
 &\text{normalizer state} \\
h_t \ &= \ o_t \ \tilde{h}_t \ ,
  &\tilde{h}_t \ &= \ c_t / n_t
&\text{hidden state} \\
z_t \ &= \ \varphi \left( \tilde{z}_t \right) \ , 
  &\tilde{z}_t \ &=  \ \mathbf{w}^\top_{z} \ \mathbf{x}_t \ + \
  r_{z}  h_{t-1} \ + \  b_{z}
  &\text{cell input} \\
i_t \ &= \ \!\exp \left( \tilde{i}_t  \right) \ , 
  &\tilde{i}_t \ &= \ \mathbf{w}^\top_{i} \ \mathbf{x}_t \ + \
  r_{i}  \ h_{t-1} \ + \  b_{i}
  &\text{input gate} \\
f_t \ &= \ \sigma \left(  \tilde{f}_t \right) \ \text{OR} \
\!\exp \left(  \tilde{f}_t \right) \ ,
  &\tilde{f}_t \ &= \ \mathbf{w}^\top_{f} \ \mathbf{x}_t  \ + \
  r_{f}  \ h_{t-1} \ + \  b_{f}
  &\text{forget gate} \\
o_t \ &= \ \sigma \left( \tilde{o}_t \right) \ , 
  &\tilde{o}_t  \ &= \ \mathbf{w}^\top_{o} \ \mathbf{x}_t \ + \
  r_{o}  \ h_{t-1} \ + \  b_{o}
  &\text{output gate} 
\end{align}

We transfer the original LSTM gating techniques, i.e., input- and/or hidden-dependent gating plus bias term, to the new architectures. Exponential activation functions can lead to large values that cause overflows. Therefore, we stabilize gates with an additional state \mbox{$m_t$~\citep{Milakov:18arxiv}}:

\begin{align}
\label{eq:slstmstabil}
m_t \ &= \ \max \left( \log ( f_t ) + m_{t-1} , \log ( i_t ) \right) &\text{stabilizer state} \\
i'_t \ &= \ \!\exp \left( \log \left ( i_t \right) - m_t \right) \ = \ \!\exp \left( \tilde{i}_t - m_t \right) \ \, 
  &\text{stabil. input gate} \\
f'_t \ &= \ \!\exp \left( \log \left( f_t \right) + m_{t-1} - m_t \right) \ \, 
  &\text{stabil. forget gate}
\end{align}

Appendix~\ref{sec:appsLSTM} demonstrates that substituting $f_t$ with $f'_t$ and $i_t$ with $i'_t$ during the forward computation does not alter either the network's overall output or the gradients of the loss relative to its parameters.

\paragraph{Novel Approach to Memory Mixing.} Similar to the conventional LSTM, sLSTM is capable of supporting several memory cells (refer to Appendix~\ref{sec:appsLSTM}). The use of multiple memory cells facilitates memory mixing, which is achieved through recurrent links $\mathbf{R}_{\mathbf{z}}$, $\mathbf{R}_{\mathbf{i}}$, $\mathbf{R}_{\mathbf{f}}$, and $\mathbf{R}_{\mathbf{o}}$ extending from the hidden state vector $\mathbf{h}$ toward the memory cell input $\mathbf{z}$ and the respective gating variables $\mathbf{i}$, $\mathbf{f}$, and $\mathbf{o}$. Exponential gating introduces a distinctive mechanism for memory mixing. The revised sLSTM architecture allows several heads, with internal memory mixing permitted within each head, but not between different heads. The combination of head-based structure and exponential gating in sLSTM forms a fundamentally new modality for memory mixing.

\subsection{mLSTM}
\label{sec:mLSTM}

To augment the storage capabilities of LSTMs, we replace the conventional memory cell, represented as a scalar $c \in \mathbb{R}$, with a matrix $\mathbf{C} \in \mathbb{R}^{d \times d}$. Consequently, retrieval operations utilize matrix multiplication. At each time step $t$, our objective is to encode a vector pair: the key $\mathbf{k}_t \in \mathbb{R}^d$ and the value $\mathbf{v}_t \in \mathbb{R}^d$—terminology borrowed from the Transformer architecture. Subsequently, at time $t + \tau$, we employ a query vector $\mathbf{q}_{t+\tau} \in \mathbb{R}^d$ to recover the value $\mathbf{v}_t$. This corresponds to the framework of Bidirectional Associative Memories (BAMs) \citep{Kohonen:72,Anderson:72,Nakano:72,Anderson:77}.

The covariance update rule~\citep{Sejnowski:77,Dayan:91} for storing a key-value pair is

\begin{align}
\mathbf{C}_t \ &= \ \mathbf{C}_{t-1} \ + \ \mathbf{v}_{t} \ \mathbf{k}_t^\top \ .
\end{align}

We postulate the application of layer normalization preceding the projection of input vectors to keys and values, ensuring these resulting vectors possess a zero mean. The covariance update mechanism is considered optimal~\citep{Dayan:91} as it achieves maximal differentiation among retrieved binary vectors, thereby enhancing the signal-to-noise ratio. Greater separability is achievable when restricting retrieval to interactions between pairs and accepting quadratic computational overhead, akin to the mechanism found in attention models~\citep{Krotov:16,Krotov:17,Ramsauer:21}. Notably, the covariance update rule mirrors the process found in Fast Weight Programmers \citep{Schmidhuber:92ncfastweights,Schlag:21}, which later incorporated a fixed decay factor for $\mathbf{C}_{t-1}$ and a constant learning coefficient applied to 
$\mathbf{v}_{t} \mathbf{k}_t ^\top$~\citep{Ba:16}. Adhering to this rationale, we embed the covariance update method within the LSTM architecture, mapping the forget gate to the decay parameter, the input gate to the learning rate, and leveraging the output gate to adjust the magnitude of the retrieved vector.

For this matrix memory, the normalizer state is the weighted sum of key vectors, where each key vector is weighted by the input gate and all future forget gates. Again, the normalizer state keeps record of the strength of the gates. Since the dot product between query and normalizer state can be close to zero, we use the absolute value of this dot product and lower bound it by a threshold (typically 1.0) as done previously~\citep{Sun:23arxiv}. The mLSTM forward pass is:

\begin{align}
\label{eq:mlstm_recurrent_begin}
\mathbf{C}_t \ &= \  f_t \ \mathbf{C}_{t-1} \ + \ 
  i_t \ \mathbf{v}_t \ \mathbf{k}_t^\top & &  &\text{cell state} \\
\mathbf{n}_t \ &= \  f_t \ \mathbf{n}_{t-1} \ + \ 
  i_t \ \mathbf{k}_t & &  &\text{normalizer state} \\
\mathbf{h}_t  \ &= \ \mathbf{o}_t \ \odot \ \tilde{\mathbf{h}}_t
\ , \qquad \qquad 
\tilde{\mathbf{h}}_t \ = \   \mathbf{C}_t \mathbf{q}_t \ / \ 
  \max \left\{ |\mathbf{n}_t^\top \mathbf{q}_t|, 1 \right\} 
   &&&\text{hidden state} \\
\mathbf{q}_t \ &= \ \mathbf{W}_q \ \mathbf{x}_t \ + \ \mathbf{b}_q  & &  &\text{query input} \\
\mathbf{k}_t \ &= \ \frac{1}{\sqrt{d}} \mathbf{W}_k \ \mathbf{x}_t \ + \ \mathbf{b}_k  & &  &\text{key input} \\
\mathbf{v}_t \ &= \ \mathbf{W}_v \ \mathbf{x}_t \ + \ \mathbf{b}_v  & &  &\text{value input} \\
i_t \ &= \ \!\exp \!  \! \left( \tilde{i}_t  \right) \ , 
 \qquad \qquad \ \ \ \ \,
  \tilde{i}_t \ = \ \mathbf{w}^\top_{i} \ \mathbf{x}_t \ + \  b_{i} \ \ \ 
  &&&\text{input gate} \\
f_t \ &= \ \sigma \!  \left(  \tilde{f}_t \right) \ \text{OR} \ \!\exp \!  \! \left(  \tilde{f}_t \right) , \ \ \: \!
\tilde{f}_t \ = \ \mathbf{w}^\top_{f} \ \mathbf{x}_t  \ + \
  b_{f}
  &&& \text{forget gate} \\
\label{eq:mlstm_recurrent_end}
\mathbf{o}_t \ &= \ \sigma \left( \tilde{\mathbf{o}}_t \right) \ , \qquad \qquad \quad \ \ \ \,
  \tilde{\mathbf{o}}_t  \ = \ \mathbf{W}_{\mathbf{o}} \ \mathbf{x}_t \ + \
  \mathbf{b}_{\mathbf{o}}
  &&&\text{output gate} 
\end{align}

mLSTM is capable of accommodating several memory cells, similar to the architecture of traditional LSTM. In the context of mLSTM, utilizing multiple heads or multiple cells yields equivalent outcomes, given the absence of memory mixing. To address the instability of mLSTM's exponential gating mechanisms, we employ stabilization methods identical to those used with sLSTM (refer to Equation~\ref{eq:slstmstabil}). 
Because mLSTM does not incorporate memory mixing, its recurrent formulation can be recast to operate in parallel. Further information is provided in Appendix~\ref{sec:appmLSTM}.

\subsection{xLSTM Architecture}
\label{sec:xLSTMarch}

\paragraph{xLSTM Blocks.}
\label{sec:xLSTMblock}
The xLSTM block aims to generate a non-linear, high-dimensional representation of prior sequence information, facilitating clearer distinction between different contextual histories. Such differentiation is crucial for accurate prediction of subsequent sequence elements (e.g., tokens). Our approach is grounded in Cover’s Theorem~\citep{Cover:65}, which posits that non-linear embeddings in higher dimensional spaces increase the probability that distinct patterns can be separated by a linear function, as compared to their separability in the original space. We evaluate two designs for residual block structures: (i) A residual block with up-projection applied after non-linear summarization—similar to the architecture used in Transformers—where the past is summarized in the original space via a non-linear operation, followed by a linear transformation into a higher dimensional space, a subsequent non-linear activation, and finally, a linear mapping back to the original dimension. This process is illustrated in the left panel of Figure~\ref{fig:Backbones} and the third column of Figure~\ref{fig:xlstm_sketch}. A comprehensive depiction is provided in Figure~\ref{fig:appsLSTM_detailed} in the appendix. (ii) Alternatively, a residual block can employ a pre up-projection—reminiscent of architectures in State Space Models—which initially linearly projects the input to a higher-dimensional space.

It summarizes previous information in a non-linear fashion within a high-dimensional space, after which it applies a linear transformation to revert to the initial feature space. When constructing an xLSTM block that incorporates an sLSTM, the post up-projection block is utilized. Alternatively, if the xLSTM block contains an mLSTM, we employ the pre up-projection block, reflecting the increase in memory capacity associated with moving into the high-dimensional space. For further illustration, please consult the left panel of Figure~\ref{fig:Backbones} and the third column of Figure~\ref{fig:xlstm_sketch}, or refer to Figure~\ref{fig:appsLSTM_detailed} in the appendix.

\paragraph{xLSTM Architecture. }
\label{sec:xLSTMarchitecure}
The xLSTM architecture is developed by stacking component blocks using residual connections, as described in~\citep{Srivastava:15,He:16}. Our implementation adopts pre-LayerNorm~\citep{Ba:16b} residual structures, aligning with those prevalent in modern Large Language Models. Refer to the final column in Figure~\ref{fig:xlstm_sketch} for details.

\subsection{Memory and Speed Considerations}

In contrast to Transformer architectures, xLSTM networks exhibit linear computational requirements and fixed memory usage relative to sequence length. The compressive nature of xLSTM memory makes it particularly advantageous for deployment in industrial contexts and for edge computing scenarios.

mLSTM’s memory does not rely on trainable parameters, yet demands considerable computational effort owing to its $d \times d$ matrix memory and associated $d \times d$ update operations. This design represents a trade-off: enhanced memory capacity comes at the expense of increased computational complexity. However, leveraging GPU parallelism mitigates most of this overhead, resulting in minimal impact on real-time execution.

The parallelization capability of mLSTM resembles that of FlashAttention~\citep{Dao:22,Dao:23} and GLA~\citep{Yang:23arxiv}, whereas sLSTM does not allow parallel execution because memory mixing (hidden-to-hidden interactions) introduces sequential dependencies. Nonetheless, we have engineered an efficient CUDA implementation, incorporating GPU memory optimizations down to the register level; this approach yields a runtime that is typically less than twice that of mLSTM.

\section{Related Work}
\label{sec:related}

\paragraph{Linear Attention.} Multiple strategies have been introduced to mitigate the Transformer’s quadratic computational burden relative to context length and achieve attention scaling that is linear in this parameter. Synthesizer generates synthetic attention scores independently of token--token relationships~\citep{Tay:20arxiv}. Linformer models self-attention using a low-rank matrix formulation, yielding a linear approximation~\citep{Wang:20arxiv}. Linear Transformer modifies the attention mechanism so it is inherently linear in computation~\citep{Katharopoulos:20}. Performer applies a positive orthogonal random features technique to produce a linearized estimation of attention’s softmax operation~\citep{Choromanski:21}. Alternatives to conventional attention, such as rapid and extensive convolutions, have been implemented in Structured Global Convolution (SGConv)~\citep{Li:22} and Hyena Hierarchy~\citep{Poli:23}.

\paragraph{State Space Models.} In recent times, State Space Models (SSMs) have gained significant attention, primarily due to their linear computational complexity with respect to context length and their strong performance relative to Transformer architectures. Initial models in this class include the Structured State Space sequence model (S4)~\citep{Gu:21}. Subsequent advancements brought forth the Diagonal State Space (DSS) model~\citep{Gupta:22}, Gated State Space (GSS) models~\citep{Mehta:22}, S5 model~\citep{Smith:22}, Bidirectional Gated SSM (BiGS)~\citep{Wang:22}, H3 model~\citep{Fu:23}, as well as Mamba~\citep{Gu:24arxiv}.

\paragraph{Recurrent Neural Networks.} Within the realm of long-context processing, Recurrent Neural Networks (RNNs) have been proposed as potential alternatives to Transformers and attention-based methods, primarily due to their linear scalability with respect to context length. Models employing Deep Linear Recurrent Units (LRUs) have achieved notable advances in language modeling~\citep{Orvieto:23, De:24arxiv}, with Hierarchically Gated Linear RNN (HGRN)~\citep{Qin:23} and its successor HGRN2~\citep{Qin:24arxiv} demonstrating similar success. Among RNN variants tailored for large-scale language modeling, the RWKV architecture~\citep{Peng:23arxivshort, Peng:24arxivshort} has emerged as a particularly strong competitor, achieving performance levels on par with Transformer-based approaches.

\paragraph{Gating.}
Gating, originally central to LSTM architectures, has been extensively adapted and re-envisioned in numerous modern models. Notable instances of gating mechanisms can be found in HGRN~\citep{Qin:23}, HGRN2~\citep{Qin:24arxiv}, Gated Linear Attention (GLA)~\citep{Yang:23arxiv}, Gated State Space (GSS) models~\citep{Mehta:22}, Bidirectional Gated SSM (BiGS)~\citep{Wang:22}, Moving Average Equipped Gated Attention (MEGA)~\citep{Ma:22}, RWKV~\citep{Peng:23arxivshort}, and Mamba~\citep{Gu:24arxiv}.

\paragraph{Covariance Update Rule.} 
In order to expand the storage capabilities of the mLSTM cell, a matrix memory utilizing a covariance update rule was introduced. Related strategies employing such updates can be seen in Fast Weight Programmers~\citep{Schmidhuber:92ncfastweights,Schlag:21}, RWKV-5 and RWKV-6~\citep{Peng:24arxivshort}, Retention~\citep{Sun:23arxiv}, Linear Transformer~\citep{Katharopoulos:20}, and HGRN2~\citep{Qin:24arxiv}.

\paragraph{Most Related.} The xLSTM bears conceptual resemblance to models such as Retention~\citep{Sun:23arxiv}, RWKV~\citep{Peng:23arxivshort,Peng:24arxivshort}, and HGRN2~\citep{Qin:24arxiv}. These architectures incorporate key elements like matrix memory and gating mechanisms. Nonetheless, unlike the newly introduced sLSTM, these models do not support memory mixing. Memory mixing is pivotal for addressing state tracking challenges, rendering LSTMs superior in expressiveness compared to State Space Models (SSMs) and Transformers~\citep{Merrill:24,Deletang:23}. Efficient state tracking is crucial when executing code or monitoring entities throughout extended narratives.

\paragraph{Residually Stacking Architectures.} The design of xLSTM architectures adheres to the standard practice of residual stacking of modular components, a strategy prevalent in most modern deep learning models~\citep{Srivastava:15,He:16}.

This approach has been foundational in enabling very deep convolutional networks~\citep{He:16} and has similarly been adopted by Transformers~\citep{Vaswani:17}. In turn, the Transformer architecture powers the development of Large Language Models (LLMs) such as GPT-3~\citep{Brown:20short}, ChatGPT~\citep{Schulman:22short}, GPT-4~\citep{Achiam:23gpt4short}, Megatron-LM~\citep{Shoeybi:19}, Gopher~\citep{Rae:21short}, ERNIE 3.0 Titan~\citep{Wang:21arxivshort}, GLaM ~\citep{Du:21glamshort}, Chinese M6~\citep{Lin:21}, AlexaTM 20B~\citep{Soltan:22}, OPT~\citep{Zhang:22opt}, Chinchilla~\citep{Hoffmann:22short}, BLOOM~\citep{Scao:22arxivshort}, GLM-130B~\citep{Zeng:22arxivshort}, LaMDA~\citep{Thoppilan:22short}, PaLM~\citep{Chowdhery:22short}, Llama~\citep{Touvron:23arxiv}, and Gemini~\citep{GeminiTeam:23arxiv,Reid:24arxiv}.

\section{Experiments}
\label{sec:Exp}

This section presents an empirical evaluation of xLSTM in comparison to established approaches, with particular emphasis on language modeling. Section~\ref{sec:ExpSyntethic} explores the distinct abilities of xLSTM by employing synthetic benchmarks. 
The subsequent Section~\ref{sec:ExpComparison} reports on the validation perplexity obtained by several state-of-the-art language modeling techniques, all trained using 15B tokens sampled from SlimPajama~\citep{Soboleva:23}. Here, we also conduct ablation analyses to dissect xLSTM's performance. In addition, we examine the scaling properties of the various models, adopting methodologies similar to \citet{Kaplan:20} and \citet{Brown:20short}.
Section~\ref{sec:ExpLanguage} 
offers an in-depth investigation of language modeling capabilities. We provide a comparative analysis between xLSTM and the highest performing models from Section~\ref{sec:ExpComparison} after training them on a dataset consisting of 300B tokens from SlimPajama~\citep{Soboleva:23}. The evaluation consists of four main components: first, we measure each method's capacity to generalize to extended contexts; second, we benchmark validation perplexity and results in downstream tasks~\citep{Sutawika:24}; third, we examine model performance across 571 domains within the PALOMA language benchmark suite~\citep{Magnusson:23arxivshort}; fourth, we revisit scaling behavior analysis, this time utilizing a dataset with 20-fold more tokens.

\label{sec:xlstm_blocks_notation}
Throughout all conducted experiments, we denote the composition of xLSTM blocks by xLSTM[$a$:$b$], where $a/b$ specifies the proportion of mLSTM-based to sLSTM-based blocks. Thus, xLSTM[7:1] indicates that among eight total blocks, seven utilize mLSTM and one employs sLSTM. Accordingly, when the model comprises 48 total blocks, this partition results in 42 mLSTM-based and 6 sLSTM-based blocks. In addition, each experiment employs pre up-projection blocks for mLSTM-based blocks and post up-projection blocks for sLSTM-based blocks.

\subsection{Synthetic Tasks and Long Range Arena}
\label{sec:ExpSyntethic}
Initially, we evaluate xLSTM's enhanced exponential gating combined with memory mixing on formal language tasks~\citep{Deletang:23}. Next, we investigate the utility of the novel matrix memory in xLSTM for the Multi-Query Associative Recall benchmark~\citep{Arora:23arxiv}. Lastly, xLSTM's capabilities in handling extended sequences are analyzed using the Long Range Arena suite~\citep{Tay:21}.

\paragraph{Evaluation of Exponential Gating with Memory Mixing in xLSTM.} We evaluate the efficacy of xLSTM’s exponential gating with memory mixing, which theoretically permits efficient state tracking as posited in recent literature~\citep{Merrill:24, Merrill:23}. Building upon the formal language experiments in \citet{Deletang:23}, we adapt these tasks to allow for training across variable sequence lengths in order to assess extrapolation capabilities. Comprehensive details regarding the suite of tasks and thorough experimental outcomes can be found in Appendix~\ref{sec:appExpSynthetic-formalLang}. Our study benchmarks xLSTM against a range of alternative approaches, including variants of Transformers, State Space Models, and classic Recurrent Neural Networks, measuring performance strictly on task-specific token predictions. Accuracy scores are normalized to lie between 0 (reflecting chance performance) and 1 (indicating flawless accuracy). Comparative analysis is carried out using two-block variants for each of the following: xLSTM[0:1] (sLSTM component alone), xLSTM[1:0] (mLSTM only), xLSTM[1:1], Llama, Mamba, RWKV, Retention, Hyena, LSTM, as well as LSTM integrated within Transformer blocks (denoted as LSTM (Block)). Figure~\ref{fig:formal-main} presents a summary of these results. Notably, architectures such as standard Transformers or State Space Models lacking memory mixing mechanisms (and thereby state tracking capacity) fail on tasks requiring regular grammar recognition, such as the parity task. These observations corroborate prior work indicating that, in terms of computational expressiveness for such tasks, Transformers and State Space Models fall short compared to RNNs~\citep{Merrill:24, Merrill:23, Deletang:23}.

\paragraph{Test of xLSTM's Memory Capacities on Associative Recall Tasks.} This study evaluates the new matrix-based memory mechanism introduced in xLSTM by assessing its ability to store and retrieve information within the Multi-Query Associative Recall task~\citep{Arora:23arxiv}. Each input sequence consists of key-value pairs, sampled at random from an extensive vocabulary, which must be remembered and accessed later in the task. To further raise the challenge from the standard formulation, we increase both the quantity of key-value pairs to 256 and the context window length up to 2048 tokens, thus providing a more rigorous assessment of the memory capacities across different models. Our comparison covers 2-block model variants, including Llama, Mamba, RWKV-5, RWKV-6, xLSTM[1:1], and xLSTM[1:0], with recall accuracy used as the evaluation criterion. Since the memory of Transformer-based architectures (e.g., Llama) scales exponentially with the embedding dimension~\citep{Ramsauer:21}, they represent the benchmark for this task. Figure~\ref{fig:mqar-main} summarizes the findings. Among non-Transformer models, xLSTM[1:1] outperforms all others, including on smaller model configurations. Notably, incorporating the sLSTM block does not impair recall capacity; rather, it improves memory performance, especially under the most demanding conditions with 256 key-value associations. Supplementary results, including extrapolation experiments suggesting that xLSTM's memory can generalize to context lengths surpassing those used during training, can be found in Appendix~\ref{sec:appExpSynthetic-mqar}.

\paragraph{Test of xLSTM's Long Context Capabilities on Long Range Arena.} In evaluating xLSTM's ability to process extended sequences and sizeable contexts, we benchmark it against alternative approaches using the Long Range Arena~\citep{Tay:21}. Across the various tasks, xLSTM reliably achieves high performance, indicating that its architecture is particularly effective for the diverse challenges associated with long-context scenarios.

For more details, see Appendix~\ref{sec:appExpSynthetic-lra}.

\subsection{Method Comparison and Ablation Study}
\label{sec:ExpComparison}

To investigate the central inquiry of this work—specifically, the potential performance of our novel LSTM variants at scale within language modelling—we conduct training of xLSTMs, Transformers, State Space Models, along with additional techniques, on 15B tokens extracted from SlimPajama, utilizing an identical auto-regressive configuration. The resulting models are evaluated using the validation set, and ablation analyses are carried out for the xLSTMs.

\paragraph{Comparison between xLSTM and Existing Approaches.} To benchmark performance, we trained models using 15B tokens sourced from SlimPajama~\citep{Soboleva:23}, and assessed their perplexity on the corresponding validation dataset. Our evaluation encompassed several approaches: xLSTM (introduced here), GPT-3 (Transformer)~\citep{Brown:20short}, Llama (Transformer)~\citep{Touvron:23arxiv}, H3 (SSM)~\citep{Fu:23}, Mamba (SSM)~\citep{Gu:24arxiv}, RWKV-4 (RNN)~\citep{Peng:23arxivshort}, RWKV-5 (RNN)~\citep{Peng:24arxivshort}, RWKV-6 (RNN)~\citep{Peng:24arxivshort}, GLA (linear Transformer)~\citep{Yang:23arxiv}, HGRN (RNN)~\citep{Qin:23}, HGRN2 (RNN)~\citep{Qin:24arxiv}, RetNet (linear Transformer)~\citep{Sun:23arxiv}, Hyena (linear Transformer)~\citep{Poli:23}, xLSTM[1:0], and xLSTM[7:1] (refer to Section~\ref{sec:xlstm_blocks_notation} for notation). Training was conducted with mixed precision; however, in the case of RWKV-5, RWKV-6, GLA, and HGRN2, the mixed precision approach omitted the use of PyTorch’s automated mixed precision capabilities (see Appendix Section~\ref{sec:appGenTrainProc} for elaboration).

The algorithms evaluated fall under three groups: (a) Transformers, (b) State Space Models (SSMs), and (c) Recurrent Neural Networks (RNNs) as well as linear Transformers. Linear Transformers utilize alternative linear mechanisms in lieu of conventional Transformer attention. Each model was constructed to align with the architectural scale of GPT-3 comprising 350M parameters (embedding dimension 1024, 24 residual layers). Notably, only GPT-3 implements shared weights between token and output embeddings, yielding a slight reduction in parameter count. As shown in Table~\ref{tab:spaj15b_model_results}, xLSTM achieves the lowest validation perplexity across all methods surveyed. Further methodological information is detailed in Appendix~\ref{sec:appExpComparison}. Figure~\ref{fig:temp_sclaw15B} illustrates the scaling properties observed during this evaluation, suggesting that xLSTM’s advantages persist with increased model capacity.

\paragraph{Ablation Studies.}
Table~\ref{tab:spaj15b_model_results} and Figure~\ref{fig:temp_sclaw15B} provide evidence that xLSTM attains strong language modeling performance when trained on 15B tokens from SlimPajama. To systematically assess the impact of each modification leading from LSTM to xLSTM, we incrementally transform a standard LSTM architecture into that of xLSTM. The process begins with embedding LSTM layers within pre-LayerNorm residual backbones. Next, the model incorporates a post up-projection block. The final adjustment involves the inclusion of exponential gating mechanisms and matrix memory.

The results are shown in Table~\ref{tab:ablstudies} (top). 

The ablation experiments link the notable enhancement in performance to the combined effects of exponential gating and matrix memory. Furthermore, recognizing the critical role of gating within RNNs and State Space Models, we systematically examine various gating strategies. The results summarized in Table~\ref{tab:ablstudies} (bottom) indicate that configuring each gate to be adaptive and responsive to the input yields progressively improved outcomes. Further analyses concerning specific backbone elements are provided in Appendix~\ref{sec:appAblDetails}.

\subsection{xLSTM as Large Language Model}
\label{sec:ExpLanguage}

To conclude this work, we conduct comprehensive experiments in large-scale language modeling, investigating xLSTM's viability as a large language model (LLM). Accordingly, we expand the training dataset, utilizing 300B tokens drawn from SlimPajama. This token quantity aligns with previous efforts such as Mamba~\citep{Gu:24arxiv} and Griffin~\citep{De:24arxiv}.

Our comparison includes xLSTM alongside RWKV-4, Llama, and Mamba, each representing a distinct methodological category described in Section~\ref{sec:ExpComparison}. RWKV-4 serves as the RNN benchmark, given that suitable training precision configurations for RWKV-5, RWKV-6, and HGRN2 were determined only after the onset of the 300B token experiments (refer to Appendix~\ref{sec:appGenTrainProc} for details).

Models are trained in various sizes: 125M, 350M, 760M, and 1.3B. All models are evaluated for their ability to extrapolate to longer sequence lengths and their validation performance. Their effectiveness is further assessed through downstream tasks, their language modeling capabilities across 571 text domains from the PALOMA benchmark, and, ultimately, their behavior according to established scaling laws.

\paragraph{Sequence Length Extrapolation.}
\label{sec:spaj300B_ctx_extrapolation}
To begin, we evaluate the sequence length extrapolation capabilities of 1.3B-parameter versions of xLSTM, RWKV-4, Llama, and Mamba. These models are trained using a fixed context length of 2048, and subsequently assessed on context lengths extending up to 16384 tokens. Results can be found in Figure~\ref{fig:spaj300B_seq_extrapolation}. Notably, xLSTM models exhibit superior performance by retaining low perplexity across extended context lengths, setting them apart from the other approaches.

\paragraph{Validation Perplexity and Downstream Tasks.}
\label{sec:spaj_downstream_eval}
Next, across all parameter scales, we measure the effectiveness of xLSTM, RWKV-4, Llama, and Mamba on the SlimPajama validation split, focusing on next token prediction, as well as on a set of downstream tasks designed to assess common sense reasoning abilities.

The third column in Table~\ref{tab:lmeval} presents the perplexities on the validation set for various methods. Among all evaluated model sizes, both xLSTM[1:0] and xLSTM[7:1] consistently achieve the lowest perplexity scores, establishing them as the top-performing models on this metric. The remaining columns in Table~\ref{tab:lmeval} display results for downstream tasks. Across nearly all tasks and model sizes, xLSTM outperforms competing methods. The only exception is the ARC task, where Mamba sometimes exhibits superior results. Additional information is available in Appendix~\ref{sec:appExpLanguage}.

\paragraph{Performance on PALOMA Language Tasks.} In the third analysis, we evaluate next token prediction accuracy across all model variants—xLSTM, RWKV-4, Llama, and Mamba—on the PALOMA language tasks~\citep{Magnusson:23arxivshort}. Performance is quantified using perplexity scores for next token prediction spanning 571 distinct textual domains, which include sources from nytimes.com as well as Reddit's r/depression.

Table~\ref{tab:ppleval} presents the perplexity results, distinguishing between broad language modeling assessments (first seven columns) and more specific domain-based benchmarks (final five columns). xLSTM[1:0] demonstrates stronger results in language tasks than xLSTM[7:1].

Specifically, xLSTM[1:0] achieves lower perplexity compared to Mamba in 568 out of 571 domains (99.5\%), outperforms Llama in 486 domains (85.1\%), and surpasses RWKV-4 in 570 domains (99.8\%). Detailed breakdowns are available in Appendix~\ref{sec:appExpLanguage}.

\paragraph{Scaling Laws.} Next, we analyze the observed power-law scaling trends, which enable predictions about model performance as model size increases~\citep{Kaplan:20,Brown:20short}.
The scaling relationships are illustrated in Figure~\ref{fig:temp_sclaw300B}. Although the models display broadly consistent scaling patterns, they are distinguished by varying baseline levels. RWKV-4 shows the weakest performance, followed by Llama and then Mamba. In comparison, xLSTM outperforms Mamba, with the performance gap between xLSTM and Mamba resembling that between Mamba and Llama. These scaling patterns suggest that xLSTM is expected to maintain its performance advantages relative to both Transformers and State-Space models as model size increases.

\textbf{Generation Times and Maximal Throughput.} In Figure~\ref{fig:inference_time_generation}, we present measurements of text generation times for various xLSTM variants at the 1.3B parameter scale. Corresponding maximal throughput results are displayed in Figure~\ref{fig:inference_time_decoding_throughput} (left). Our evaluation includes similarly-sized models implemented in Mamba, Llama, and RWKV sourced from HuggingFace, with the Llama model incorporating a static key-value cache. During the experiments, it was not possible to compile either a complete key-value cache for the Transformer Model or the Mamba model with \texttt{torch.compile}. All text generation tests were conducted with batch size set to 1 and a pre-fill value of 16, a configuration particularly advantageous for the Transformer. 

As observed in Figure~\ref{fig:inference_time_generation}, both xLSTM and other recurrent architectures such as Mamba and RWKV-4 demonstrate linear complexity in generation times, in contrast to the quadratic scaling characteristic of Llama. For the decoding throughput comparison, we examine performance across different batch sizes and pre-fill settings for the Llama baseline. As depicted in Figure~\ref{fig:inference_time_decoding_throughput} (right), xLSTM supports considerably larger batch sizes than Llama, attributed to its constant memory requirements, and, as a consequence, delivers the highest throughput among the tested models.

\section{Limitations}

(i) Unlike mLSTM, the way sLSTM mixes memory impedes operations that can be parallelized, making rapid parallel execution unattainable. Despite this, we have designed an efficient CUDA kernel for sLSTM, which currently operates at less than twice the runtime of our parallel mLSTM implementation. (ii) The existing CUDA kernels for mLSTM lack optimization, resulting in performance that is approximately four times slower than FlashAttention or the scan routine used in Mamba. Optimized CUDA kernels, similar to the approach used by FlashAttention, could substantially speed up mLSTM computations. (iii) The matrix-based memory in mLSTM incurs significant computational complexity since it processes $d \times d$ matrices. However, both memory update and access routines are parameter-free and benefit from parallel execution via standard matrix operations, which keeps the added computational burden minimal in terms of actual runtime. (iv) Careful consideration must be given when initializing the forget gates. (v) Because matrix memory is not tied to sequence length, longer sequences could potentially overwhelm memory during extended contexts. However, this has not proven problematic for contexts up to 16k tokens, as described in Section~\ref{sec:spaj300B_ctx_extrapolation}. (vi) Given the heavy computational demands of large language model experiments, neither the architecture nor the hyperparameter selection underwent complete optimization—particularly for larger xLSTM designs. We expect that achieving optimal results with xLSTM will require a thorough and systematic optimization process.

\section{Conclusion}
We have partially addressed our central inquiry: What advancements can be achieved in language modeling by scaling LSTM architectures to the billion-parameter regime? Our findings suggest that these scaled models can perform at least as well as existing solutions such as Transformers and State Space Models. We introduced xLSTM, a variation of LSTM incorporating exponential gating mechanisms with memory mixing and a reconfigured memory architecture. Empirically, xLSTM models demonstrate competitive or superior performance on language modeling tasks relative to state-of-the-art models including Transformers and State Space Models. Analysis of scaling behaviors suggests that increasing the size of xLSTM models positions them as strong challengers to contemporary Large Language Models constructed with Transformer frameworks. Furthermore, xLSTM holds promise for transformative applications in areas such as Reinforcement Learning, Time Series Prediction, and the modeling of physical processes.

\end{document}